\subsubsection{Matrix Packaging of the Dirac Pair}
\label{sec:matrix-dirac-packaging}
The intrinsic ordered-pair equations in the wave-mechanics section can be
packaged by an auxiliary matrix function. This packaging translates those
equations without introducing a new physical assumption.

\paragraph{Auxiliary block map.}
The block identity needed below is defined by
\begin{equation}
\I:=
\begin{bmatrix}
1&0\\
0&1
\end{bmatrix}.
\label{eq:W-block-identity}
\end{equation}
The two-input block-diagonal function needed below is defined by
\begin{equation}
\eqfunc{diag}(\pv{A},\pv{B}):=
\begin{bmatrix}
\pv{A}&0\\
0&\pv{B}
\end{bmatrix}.
\label{eq:W-block-diag}
\end{equation}
A paravector \(\pv{Z}\in\CPspace\), including any member of
\(\RPspace\), is packaged into the two-by-two block map
\begin{equation}
\Wmat(\pv{Z}):=
\begin{bmatrix}
0&\pv{Z}\\
\adjf{\pv{Z}}&0
\end{bmatrix}.
\label{eq:W-map}
\end{equation}
For \(a,b\in\R\), its complex linearity is
\begin{equation}
\Wmat\bigl((a+ib)\pv{Z}\bigr)
=(a+ib)\,\Wmat(\pv{Z}).
\label{eq:W-linearity}
\end{equation}
Its square is
\begin{equation}
\Wmat(\pv{Z})^2
=\eqfunc{diag}(\detf{\pv{Z}},\detf{\pv{Z}})
=\detf{\pv{Z}}\,\I.
\label{eq:W-square}
\end{equation}
The block map \(\Wmat(\pv{Z})\) is a two-by-two matrix over the paravector
algebra. After each block is replaced by its Pauli matrix, it is a
four-by-four complex matrix. Addition and multiplication by ordinary
scalars are preserved, so the packaging map is linear. General products are
not preserved, so it is not an additional representation of the algebra.
Equation \eqref{eq:W-square} is obtained because its
diagonal block products are reduced to the determinant in
Eq.\ \eqref{eq:minkowski-form}; the off-diagonal blocks are zero
\cite{hestenes1966,doran2003}.

If \(\detf{\pv{Z}}\neq0\), Eq.\ \eqref{eq:W-square} shows that the powers of
the block matrix are governed by the scalar equation
\(x^2-\detf{\pv{Z}}=0\). Its two roots are distinct; selecting one square
root fixes their labels. An invertible change of basis
can then put the block matrix into diagonal form; this property is called
diagonalizability. If
\(\detf{\pv{Z}}=0\) and \(\pv{Z}\neq0\), then
\(\Wmat(\pv{Z})^2=0\) while \(\Wmat(\pv{Z})\neq0\). A nonzero matrix whose
positive power is zero is called nilpotent; such a matrix cannot be put
into diagonal form. For blocks containing differential operations or other
quantities whose order cannot be exchanged, the unreduced statement is
\[
\Wmat(\pv{Z})^2
=\eqfunc{diag}(\pv{Z}\cherm{\pv{Z}},\cherm{\pv{Z}}\pv{Z}).
\]
A separate operator calculation is required before either ordered product can
be called a determinant or both can be reduced to one scalar multiple of the
identity.

For a real normalized paravector \(\pv{U}\) with
\(\detf{\pv{U}}=1\), Eq.\ \eqref{eq:W-square} gives
\(\Wmat(\pv{U})^2=\I\). The two auxiliary projector matrices are therefore
\begin{equation}
\frac12\bigl(\I\pm\Wmat(\pv{U})\bigr).
\label{eq:W-projectors}
\end{equation}
No additional projector function is introduced by this optional block
construction. Direct multiplication gives the complete algebraic properties:
\begin{equation}
\begin{aligned}
\left(\frac12\bigl(\I\pm\Wmat(\pv{U})\bigr)\right)^2
&=\frac12\bigl(\I\pm\Wmat(\pv{U})\bigr),\\
\frac12\bigl(\I+\Wmat(\pv{U})\bigr)
\frac12\bigl(\I-\Wmat(\pv{U})\bigr)&=0,\\
\frac12\bigl(\I+\Wmat(\pv{U})\bigr)
+\frac12\bigl(\I-\Wmat(\pv{U})\bigr)&=\I,\\
\Wmat(\pv{U})\frac12\bigl(\I\pm\Wmat(\pv{U})\bigr)
&=\pm\frac12\bigl(\I\pm\Wmat(\pv{U})\bigr).
\end{aligned}
\label{eq:W-projector-properties}
\end{equation}
For any conformable two-entry column
\(\bPsi_{\eqtext{auxiliary}}\) with paravector entries, the normalized
auxiliary condition is equivalently written as
\begin{align}
\bigl(\Wmat(\pv{U})-\I\bigr)\bPsi_{\eqtext{auxiliary}}
&=0,\notag\\
\bPsi_{\eqtext{auxiliary}}
&=\frac12\bigl(\I+\Wmat(\pv{U})\bigr)
\bPsi_{\eqtext{auxiliary}}.
\label{eq:W-normalized-condition}
\end{align}
This is only the \(+1\) eigenvalue restriction of the auxiliary map.
No particle-versus-antiparticle meaning follows from the two signs without
additional physical structure.

\paragraph{Full-matrix and entrywise Hermitian operations.}
Two different levels of operation must be distinguished for a matrix whose
entries are themselves paravectors. By separate application of the intrinsic
Hermitian operation to each entry of \(\Wmat(\pv{Z})\), the following matrix
is obtained:
\begin{equation*}
\begin{aligned}
&\eqtext{entrywise }\Hop:\qquad
\begin{bmatrix}
0&\pv{Z}\\
\adjf{\pv{Z}}&0
\end{bmatrix}\\
&\qquad\longmapsto\qquad
\begin{bmatrix}
0&\herm{\pv{Z}}\\
\conj{\pv{Z}}&0
\end{bmatrix}
=\Wmat(\herm{\pv{Z}}).
\end{aligned}
\end{equation*}
The Hermitian operation on the complete matrix also exchanges the two
off-diagonal positions. The following result is therefore obtained:
\begin{equation*}
\herm{\Wmat(\pv{Z})}
=
\begin{bmatrix}
0&\conj{\pv{Z}}\\
\herm{\pv{Z}}&0
\end{bmatrix}
=\Wmat(\conj{\pv{Z}}).
\end{equation*}
Thus entrywise application of \(\Hop\) is not the Hermitian operation of the
complete block matrix. The same final matrix as the full-matrix Hermitian
operation is obtained by entrywise intrinsic star:
\begin{equation*}
\conj{\Wmat(\pv{Z})}=\Wmat(\conj{\pv{Z}}).
\end{equation*}
In the entrywise operation, however, the block positions are not exchanged.

\paragraph{Free Dirac packaging.}
The following relations are obtained from complex linearity of the auxiliary
function:
\begin{align*}
\Wmat(i\pv{\partial})
&=i\,
\begin{bmatrix}
0&\pv{\partial}\\
\conj{\pv{\partial}}&0
\end{bmatrix}
=i\,\Wmat(\pv{\partial}),\\
\Wmat(i\pv{\partial})^2
&=-\eqfunc{diag}\bigl(
\pv{\partial}\conj{\pv{\partial}},
\conj{\pv{\partial}}\pv{\partial}
\bigr)\notag\\
&=-\Box\I.
\end{align*}
The two mass-sign factors consequently satisfy
\begin{align*}
(\Wmat(i\pv{\partial})+m\I)
(\Wmat(i\pv{\partial})-m\I)
&=-(\Box+m^2)\I,\\
(\Wmat(i\pv{\partial})-m\I)
(\Wmat(i\pv{\partial})+m\I)
&=-(\Box+m^2)\I.
\end{align*}
The free pair
\(\bPsi_{\eqtext{Dirac}}
:=(\bPsi_{\eqtext{Dirac},1},\bPsi_{\eqtext{Dirac},2})\)
can therefore be packaged as
\begin{equation*}
(\Wmat(i\pv{\partial})-m\I)
\begin{bmatrix}
\bPsi_{\eqtext{Dirac},1}\\
\bPsi_{\eqtext{Dirac},2}
\end{bmatrix}
=0.
\end{equation*}
Expansion of the two matrix rows recovers the intrinsic equations
\begin{align*}
i\pv{\partial}\bPsi_{\eqtext{Dirac},2}
&=m\bPsi_{\eqtext{Dirac},1},\\
i\conj{\pv{\partial}}\bPsi_{\eqtext{Dirac},1}
&=m\bPsi_{\eqtext{Dirac},2}.
\end{align*}

For the fixed projector used in the intrinsic derivation, the two
paravectors are represented by the columns
\begin{equation*}
\begin{aligned}
\bPsi_{\eqtext{Dirac},1}&\quad\eqtext{is represented by}\quad
\begin{bmatrix}a\\b\end{bmatrix},\\
\bPsi_{\eqtext{Dirac},2}&\quad\eqtext{is represented by}\quad
\begin{bmatrix}c\\d\end{bmatrix}.
\end{aligned}
\end{equation*}
For a free plane wave, the coupled matrix equations are
\begin{equation*}
(E-\vct{p}\cdot\sigv)
\begin{bmatrix}c\\d\end{bmatrix}
=m\begin{bmatrix}a\\b\end{bmatrix},
\qquad
(E+\vct{p}\cdot\sigv)
\begin{bmatrix}a\\b\end{bmatrix}
=m\begin{bmatrix}c\\d\end{bmatrix}.
\end{equation*}
Their four scalar entries are exactly the component equations displayed in
the intrinsic derivation. If both columns of each unrestricted Pauli matrix
are retained, two independent copies of this Dirac system are obtained.

\paragraph{Matrix form of the Klein--Gordon lift.}
The intrinsic lift in Eq.\ \eqref{eq:dirac-klein-gordon-lift} is packaged by
\begin{equation*}
\begin{bmatrix}
\bPsi_{\eqtext{Dirac},1}\\
\bPsi_{\eqtext{Dirac},2}
\end{bmatrix}
=\frac1m
\bigl(\Wmat(i\pv{\partial})+m\I\bigr)
\begin{bmatrix}
\psi_{\eqtext{KG}}\pv{a}\\
\psi_{\eqtext{KG}}\pv{b}
\end{bmatrix}.
\end{equation*}
The following result is then obtained from the free factorization:
\begin{equation*}
\bigl(\Wmat(i\pv{\partial})-m\I\bigr)
\begin{bmatrix}
\bPsi_{\eqtext{Dirac},1}\\
\bPsi_{\eqtext{Dirac},2}
\end{bmatrix}
=-\frac1m(\Box+m^2)\psi_{\eqtext{KG}}
\begin{bmatrix}
\pv{a}\\
\pv{b}
\end{bmatrix}
=0.
\end{equation*}
For \(\pv{a}=\pv{\Pi}(+)\) and \(\pv{b}=0\), the result is
\begin{equation*}
\begin{bmatrix}
\bPsi_{\eqtext{Dirac},1}\\
\bPsi_{\eqtext{Dirac},2}
\end{bmatrix}
=
\begin{bmatrix}
\psi_{\eqtext{KG}}\pv{\Pi}(+)\\
\dfrac{i}{m}(\conj{\pv{\partial}}\psi_{\eqtext{KG}})
\pv{\Pi}(+)
\end{bmatrix}.
\end{equation*}

\paragraph{Potential-coupled packaging.}
The potential-coupled pair is denoted by
\begin{equation*}
\bPsi_{\eqtext{Dirac\_with\_potential}}
:=\bigl(
\bPsi_{\eqtext{Dirac\_with\_potential},1},
\bPsi_{\eqtext{Dirac\_with\_potential},2}
\bigr).
\end{equation*}
Its field equation is represented by the block operator
\begin{equation}
\Dmat:=
\Wmat(i\pv{\partial}-q\conj{\pv{\Phi}})-m\I
=i\,\Wmat(\pv{\partial})
-q\,\Wmat(\conj{\pv{\Phi}})-m\I.
\label{eq:matrix-dirac-pair-operator}
\end{equation}
The following elementary block products are obtained:
\begin{align*}
\Wmat(\pv{\partial})^2&=\Box\I,\\
\Wmat(\pv{\partial})\Wmat(\conj{\pv{\Phi}})
&=
\begin{bmatrix}
0&\pv{\partial}\\
\conj{\pv{\partial}}&0
\end{bmatrix}
\begin{bmatrix}
0&\conj{\pv{\Phi}}\\
\pv{\Phi}&0
\end{bmatrix}\notag\\
&=\eqfunc{diag}\bigl(
\pv{\partial}\pv{\Phi},
\conj{\pv{\partial}}\conj{\pv{\Phi}}
\bigr),\\
\Wmat(\conj{\pv{\Phi}})\Wmat(\pv{\partial})
&=
\begin{bmatrix}
0&\conj{\pv{\Phi}}\\
\pv{\Phi}&0
\end{bmatrix}
\begin{bmatrix}
0&\pv{\partial}\\
\conj{\pv{\partial}}&0
\end{bmatrix}\notag\\
&=\eqfunc{diag}\bigl(
\conj{\pv{\Phi}}\conj{\pv{\partial}},
\pv{\Phi}\pv{\partial}
\bigr),\\
\Wmat(\conj{\pv{\Phi}})^2
&=
\begin{bmatrix}
0&\conj{\pv{\Phi}}\\
\pv{\Phi}&0
\end{bmatrix}^{2}\notag\\
&=\eqfunc{diag}\bigl(
\conj{\pv{\Phi}}\pv{\Phi},
\pv{\Phi}\conj{\pv{\Phi}}
\bigr)\notag\\
&=\detf{\pv{\Phi}}\I.
\end{align*}
The complete square is expanded as
\begin{align*}
\Wmat(i\pv{\partial}-q\conj{\pv{\Phi}})^2
&=-\Wmat(\pv{\partial})^2\notag\\
&\quad-iq\bigl(
\Wmat(\pv{\partial})\Wmat(\conj{\pv{\Phi}})
+\Wmat(\conj{\pv{\Phi}})\Wmat(\pv{\partial})
\bigr)\notag\\
&\quad+q^2\,\Wmat(\conj{\pv{\Phi}})^2\notag\\
&=-\eqfunc{diag}\!\Bigl(
(\pv{\partial}+iq\conj{\pv{\Phi}})
(\conj{\pv{\partial}}+iq\pv{\Phi}),\notag\\
&\qquad\qquad
(\conj{\pv{\partial}}+iq\pv{\Phi})
(\pv{\partial}+iq\conj{\pv{\Phi}})
\Bigr).
\end{align*}
The coupled field equation is therefore packaged as
\begin{equation*}
\Dmat
\begin{bmatrix}
\bPsi_{\eqtext{Dirac\_with\_potential},1}\\
\bPsi_{\eqtext{Dirac\_with\_potential},2}
\end{bmatrix}
=0.
\end{equation*}

The reversed-mass partner is
\begin{equation}
\Dmat':=
\Wmat(i\pv{\partial}-q\conj{\pv{\Phi}})+m\I.
\label{eq:matrix-dirac-pair-operator-partner}
\end{equation}
Its product with \(\Dmat\) is
\begin{align*}
\Dmat'\Dmat
&=\Wmat(i\pv{\partial}-q\conj{\pv{\Phi}})^2-m^2\I\notag\\
&=-\eqfunc{diag}\!\Bigl(
(\pv{\partial}+iq\conj{\pv{\Phi}})
(\conj{\pv{\partial}}+iq\pv{\Phi})+m^2,\notag\\
&\qquad\qquad
(\conj{\pv{\partial}}+iq\pv{\Phi})
(\pv{\partial}+iq\conj{\pv{\Phi}})+m^2
\Bigr).
\end{align*}
The two diagonal entries are precisely the two ordered second-order
paravector equations obtained from
\begin{equation*}
\mathcal D'\bigl(
\pv{\Phi},
\mathcal D(\pv{\Phi},\bPsi_{\eqtext{Dirac\_with\_potential}})
\bigr).
\end{equation*}

\paragraph{Gauge transform in the matrix packaging.}
With
\begin{align*}
\bPsi_{\eqtext{Dirac\_with\_potential}}
&=\exp(iq\lambda)\bPsi_{\eqtext{Dirac\_with\_potential}}',\\
\pv{\Phi}'&=\pv{\Phi}+\conj{\pv{\partial}}\lambda,
\end{align*}
the packaged gauge calculation is
\begin{align*}
0
&=(\Wmat(i\pv{\partial}-q\conj{\pv{\Phi}})-m\I)
\bPsi_{\eqtext{Dirac\_with\_potential}}\notag\\
&=(\Wmat(i\pv{\partial}-q\conj{\pv{\Phi}})-m\I)
\exp(iq\lambda)\bPsi_{\eqtext{Dirac\_with\_potential}}'\notag\\
&=\exp(iq\lambda)
\bigl(
\Wmat(i\pv{\partial}-q\conj{\pv{\Phi}'})-m\I
\bigr)\bPsi_{\eqtext{Dirac\_with\_potential}}'.
\end{align*}
The transformed matrix equation is therefore
\begin{equation*}
\bigl(
\Wmat(i\pv{\partial}-q\conj{\pv{\Phi}'})-m\I
\bigr)\bPsi_{\eqtext{Dirac\_with\_potential}}'=0.
\end{equation*}
